\documentclass[11pt]{article}
\usepackage[margin=25mm]{geometry}
\usepackage{fontspec}
\setmainfont{Latin Modern Roman}
\usepackage{amsmath,amssymb,mathtools}
\usepackage{unicode-math}
\setmathfont{Latin Modern Math}
\usepackage{xurl,hyperref}
\hypersetup{colorlinks=true,urlcolor=blue}
\setcounter{secnumdepth}{0}
\setlength{\parindent}{0pt}
\setlength{\parskip}{5pt}
\setlength{\emergencystretch}{3em}
\newcommand{\sref}[2]{\hyperlink{source:#1:#2}{[S#2]}}
\newcommand{\cataloguescope}[1]{\paragraph{Recorded target.}#1}
\newcommand{\R}{\mathbb R}
\newcommand{\C}{\mathbb C}
\newcommand{\N}{\mathbb N}
\newcommand{\Z}{\mathbb Z}
\newcommand{\Q}{\mathbb Q}
\newcommand{\F}{\mathbb F}
\newcommand{\eps}{\varepsilon}
\begin{document}
\section{Knots–Quivers Correspondence for All Knots}\label{statement}

\subsection{Definitions and mathematical statement}

\paragraph{Shared notation.}$\mathbb N=\{1,2,\ldots\}$; the standard number fields have their usual meanings. All additional conventions are stated locally.


For a knot $K$, colored HOMFLY--PT invariants are assembled into a generating series $P_K(x,a,q)$. A quiver is a finite directed graph, encoded by a matrix $C$; its motivic generating function $P_C$ packages its representation-theoretic invariants. The knots--quivers correspondence identifies these series after the source's grading substitutions.

\paragraph{Catalogue wording retained for traceability.}
Does the knots–quivers correspondence hold for every knot, relating its colored HOMFLY–PT generating series $P(x,a,q)$ to the motivic generating series $P_C(x,q)$ of a corresponding quiver?

\paragraph{Scope and source qualification.}The substitutions and framing conventions are part of the correspondence. Equality of one uncolored polynomial would not establish the stated all-color generating-series claim.
\subsection{Short English statement}
Determine whether every knot has a quiver whose motivic series reproduces all its colored HOMFLY--PT invariants.
\subsection{Research attempt}
\paragraph{Model, date and investigated scope.}
OpenAI GPT-6 Astra (configured model \texttt{gpt-6-astra}), ultra reasoning; generated 11 October 2026. Computed low-degree quiver coefficients and exhibited identical linear data with distinct quadratic interactions; no all-color knot correspondence.

\paragraph{Sources and literature check.}
Read the supplied definition, catalogue wording, source qualifications and reference entries for this question. The historical source/context and inherited status are retained; no fresh literature or novelty certification is claimed. No external theorem is used without proof in this attempt; explicit model assumptions are distinguished from facts about the source model.

\paragraph{What low colors can fail to determine.} In an auxiliary standard symmetric-quiver convention write
\[
 P_C(x_1,\ldots,x_m;q)=\sum_{d_i\ge0}
 \frac{(-q)^{d^TCd}\prod_i x_i^{d_i}}{\prod_i\prod_{r=1}^{d_i}(1-q^{2r})},
\]
with symmetric integer matrix \(C\) and formal \(q\). The coefficient of \(x_i\) is \((-q)^{C_{ii}}/(1-q^2)\), whereas for \(i\ne j\) the coefficient of \(x_ix_j\) is
\[
 \frac{(-q)^{C_{ii}+C_{jj}+2C_{ij}}}{(1-q^2)^2}.
\]
Thus the linear data see the diagonal entries but not off-diagonal arrows. With two vertices and both diagonal entries zero, matrices having \(C_{12}=0\) or \(C_{12}=1\) agree in every linear coefficient and disagree in the mixed quadratic coefficient by a factor \(q^2\).
\paragraph{Attempt and audit.} Matching an uncolored knot invariant to the first generating coefficient therefore leaves interaction data undetermined even before grading substitutions or framing are addressed. The calculation does not assert that either example is the quiver of a particular knot, nor that this auxiliary normalization is the only source convention.
\paragraph{Full-scope gap.} The first missing step is a construction from arbitrary knot data that supplies all quiver entries and the exact grading/framing substitution, followed by an identity of the entire colored series. Finite coefficient matches cannot exclude disagreement at a later color, and this explicit degree-two ambiguity shows where a first-coefficient reconstruction fails. No all-knot correspondence or new knot case is proved.

\paragraph{Verification and adversarial audit.}
Checked the quadratic exponent including the factor two on off-diagonal entries and kept the chosen generating convention separate from unverified source normalization. This is a GPT-6 Astra Ultra self-audit, not independent human verification. Prompt A drafting followed by prompt B scope/dependency/boundary self-audit. Checked the quadratic exponent including the factor two on off-diagonal entries and kept the chosen generating convention separate from unverified source normalization. The remaining gap is: Construct the quiver and exact grading/framing substitutions for arbitrary knots and prove the full formal series identity. Unresolved status and zero original-target completion are retained.

Coordinator checked the linear and mixed quadratic quiver coefficients and the off-diagonal ambiguity. This additional check was model-assisted, not independent human verification.

\paragraph{Tools, process and contributions.}
Python 3 and local file-reading/writing tools for manuscript serialization and structural checks; mathematical calculations were symbolic reasoning, with no numerical experiment or theorem prover. Contribution prompt A was followed by prompt B. GPT-6 Astra with ultra reasoning drafted the research and performed its stated self-audit; any additional coordinator contribution is identified in the verification and audit. A GPT-6-based Codex coordinator assembled the manuscript and checked repository integration; its exact variant and reasoning configuration were not exposed. No human mathematical contribution, independent human proof check or formal proof certification is supplied. Complete derivations and the consulted source references are included here; private reasoning transcripts are not published.

\paragraph{Final status and first remaining gap.}
\textbf{UNRESOLVED}. Construct the quiver and exact grading/framing substitutions for arbitrary knots and prove the full formal series identity.

\section{Completion Estimate}
\noindent\textbf{COMPLETION: 0\%}
The proved calculation is an elementary auxiliary result or a test of a restricted strategy; no new portion of the full source target is established.

\subsection{Sources}
\begin{itemize}
\item[\hypertarget{source:statement:1}{S1}] ULAM.AI and the UnsolvedMath Contributors, UnsolvedMath, catalogue entry. Catalogue curation under CC BY 4.0; definitions and English exposition added locally.
\url{https://huggingface.co/datasets/ulamai/UnsolvedMath}.
\item[\hypertarget{source:statement:2}{S2}] Mini-Workshop: Quantization of Complex Symplectic Varieties (2022), Oberwolfach Reports, DOI 10.4171/owr/2022/45, pp. 2629--2658. Original report and mathematical context.
\url{https://doi.org/10.4171/owr/2022/45}.
\end{itemize}
\label{end:statement}
\end{document}
